\documentclass{article}
\usepackage{/home/user1/code/tex/sty}
\usepackage{amsfonts}
\usepackage{amsmath}
\usepackage{geometry}
\usepackage{graphicx}
\usepackage{hyperref}
\usepackage{floatrow}
\newcommand{\fig}[3]{
\begin{figure}[H]
\begin{center}
\includegraphics[height=#1]{#2}
\caption{#3}
\label{#2}
\end{center}
\end{figure}}
\usepackage[backend=bibtex]{biblatex}
\addbibresource{bib.bib}
\setlength{\parindent}{0pt}
\geometry{top=1in,bottom=1in,left=1in,right=1in}
\n{thcab}{\theta_{CAB}}
\n{thabc}{\theta_{ABC}}
\begin{document}
\begin{center}
\textbf{\Large{Problem Set 5 Corrections}}\\~\\
\begin{tabular}{l l}
Student&Agustin Romero\\
Email&ar1@stanford.edu\\
Instructor&Caterina Vernieri\\
Course&Physics 252\\
Institution&Stanford University\\
Date&\today\\
\end{tabular}
\end{center}
\section*{Problem 1.a.} 
The student solution was not correct. Each quark has two possible states, $m_s=\pm \fr{1}{2}$, so there are $2\times 2\times 2=8$ possible initial states. Afterwards, there is one state with spin $\fr{3}{2}$, with 4 possible $m_s$ states, one state with spin $\fr{1}{2}$ with two possible $m_s$ states, and one other state with spin $\fr{1}{2}$ also with two possible $m_s$ states. So the final state has 8 possible spin states.
\section*{Problem 1.b.} 
The student solution was correct.
\section*{Problem 2}
The student solution was correct.
\section*{Problem 3}
The student solution was not correct. The solution said $l=1$ or $l=2$ were possible, but the solution also required the conservation of parity, which is
\e{P(\Delta)=1}
Which when compared to the parity of the final state
\e{P(p)P(\pi)(-1)^{l}=(-1)^{l+1}}
results in $l=1$ being the only allowed value.
\section*{Problem 4.a.}
The student solution was correct.
\section*{Problem 4.b.}
The student solution was correct.
\section*{Problem 4.c.}
The student solution was not correct. Because there are no isotopes of He or Li with atomic weight 4, the $\alpha$ particle must be in an $I=0$ state, not a $\ket{2,0}$ state. The hypothetical process would then be
\e{\ket{0,0}\rightarrow \ket{1,0}}
which violates isospin conservation. 
\end{document}
